\documentclass[12pt,a4paper]{article}

% --- Packages ---
\usepackage[utf8]{inputenc}
\usepackage[T1]{fontenc}
\usepackage[french]{babel}
\usepackage{geometry}
\geometry{top=2.5cm, bottom=2.5cm, left=2.5cm, right=2.5cm}
\usepackage{graphicx}
\usepackage{hyperref}
\usepackage{xcolor}
\usepackage{listings}
\usepackage{titlesec}
\usepackage{enumitem}
\usepackage{booktabs}
\usepackage{tcolorbox}

\hypersetup{
    colorlinks=true,
    linkcolor=blue,
    filecolor=magenta,
    urlcolor=cyan,
    pdftitle={Plan de Simulation GANDAL PaaS sur AWS - v3},
}

% --- Affichage du code ---
\definecolor{codegreen}{rgb}{0,0.6,0}
\definecolor{codegray}{rgb}{0.5,0.5,0.5}
\definecolor{codepurple}{rgb}{0.58,0,0.82}
\definecolor{backcolour}{rgb}{0.95,0.95,0.92}

% Le langage YAML n'est pas fourni par listings : definition necessaire
\lstdefinelanguage{yaml}{
  keywords={true,false,null,y,n},
  sensitive=false,
  comment=[l]{\#},
  morestring=[b]',
  morestring=[b]"
}

\lstdefinestyle{mystyle}{
    backgroundcolor=\color{backcolour},
    commentstyle=\color{codegreen},
    keywordstyle=\color{magenta},
    numberstyle=\tiny\color{codegray},
    stringstyle=\color{codepurple},
    basicstyle=\ttfamily\footnotesize,
    breakatwhitespace=false,
    breaklines=true,
    captionpos=b,
    keepspaces=true,
    numbers=left,
    numbersep=5pt,
    showspaces=false,
    showstringspaces=false,
    showtabs=false,
    tabsize=2
}
\lstset{style=mystyle}
% Accents francais dans les listings (compatible pdflatex)
\lstset{literate=
  {é}{{\'e}}1 {è}{{\`e}}1 {ê}{{\^e}}1 {ë}{{\"e}}1
  {à}{{\`a}}1 {â}{{\^a}}1 {î}{{\^i}}1 {ï}{{\"i}}1
  {ô}{{\^o}}1 {ù}{{\`u}}1 {û}{{\^u}}1 {ç}{{\c c}}1
  {É}{{\'E}}1 {È}{{\`E}}1 {À}{{\`A}}1 {œ}{{\oe}}1
}

% --- Titres ---
\titleformat{\section}{\Large\bfseries\color{blue!80!black}}{\thesection}{1em}{}
\titleformat{\subsection}{\large\bfseries\color{blue!60!black}}{\thesubsection}{1em}{}

% --- Page de garde ---
\title{
    \vspace{2cm}
    \textbf{\Huge Projet GANDAL} \\
    \vspace{0.5cm}
    \Large{Procédure d'implémentation de la simulation PaaS} \\
    \vspace{1cm}
    \large{Approche « outils standards » pour une portabilité maximale vers XCP-ng} \\
    \vspace{0.5cm}
    \normalsize{Version 3}
}
\author{\textbf{Cellule PaaS}}
\date{\today}

\begin{document}

\maketitle
\thispagestyle{empty}
\newpage

\tableofcontents
\newpage

% ==========================================
\section{Introduction et philosophie de la simulation}

\subsection{Objectif}
Ce document décrit la procédure pour déployer sur une VM AWS un mini-datacenter fonctionnel : une grosse VM, un cluster Kubernetes mono-nœud, et un script PaaS qui crée à la demande les conteneurs, les pods, le rattachement à un utilisateur et un nom DNS d'accès. L'objectif est de valider la mécanique Kubernetes (Namespaces, Deployments, Services, Ingress, Quotas, NetworkPolicies) sans dépendre de l'infrastructure physique finale.

\subsection{Règle d'or : zéro service managé}
AWS est utilisé uniquement comme fournisseur de machines virtuelles. Aucun service managé n'est employé (ni EKS, ni ECR, ni Route 53). Les équivalents open source sont installés à la main :
\begin{itemize}
    \item \textbf{Kubernetes} : K3s sur une instance EC2 ;
    \item \textbf{Registry} : \nolinkurl{registry:2} déployé dans le cluster ;
    \item \textbf{DNS} : enregistrement wildcard vers l'IP de la VM (voir Phase~6).
\end{itemize}

\subsection{Docker et Kubernetes : qui fait quoi}
Docker sert à \textbf{construire et publier} les images. Kubernetes (via le runtime \textit{containerd} intégré à K3s) \textbf{exécute} les conteneurs. Les images construites avec Docker sont exécutées par K3s sans modification.

\subsection{Fichiers livrés}
La procédure repose sur trois fichiers, à conserver dans un dépôt Git plutôt que dans ce document :
\begin{center}
\begin{tabular}{ll}
\toprule
\textbf{Fichier} & \textbf{Rôle} \\
\midrule
\texttt{install.sh} & Installe Docker, K3s et le registry (autonome, rejouable) \\
\texttt{registry.yaml} & Manifeste du registry (déjà inclus dans \texttt{install.sh}) \\
\texttt{gandal\_deploy.py} & Script PaaS : \texttt{deploy}, \texttt{list}, \texttt{status}, \texttt{delete} \\
\bottomrule
\end{tabular}
\end{center}

\subsection{Principes pour réduire l'effort de migration vers GANDAL}
\begin{itemize}
    \item \textbf{Le script est la source de vérité} : tout ce qui est créé pour un utilisateur (namespace, quota, NetworkPolicy, déploiement, service, ingress) l'est par \texttt{gandal\_deploy.py}, pas par des YAML appliqués à la main.
    \item \textbf{Aucune valeur dépendante de l'environnement dans le code} : domaine, classe d'Ingress et quotas sont lus dans des variables d'environnement.
    \item \textbf{Installation rejouable} : \texttt{install.sh} installe tout le socle en une commande.
\end{itemize}

\begin{tcolorbox}[colback=green!5!white,colframe=green!75!black,title=Portabilité]
La logique Kubernetes et le script Python sont réutilisables tels quels sur XCP-ng. Ce qui change est listé explicitement à la Phase~9 (adresse IP, domaine, TLS, accès Internet, éventuellement le CNI).
\end{tcolorbox}

% ==========================================
\section{Phase 1 : VM sur AWS (le « faux » XCP-ng)}

\subsection{But}
Créer la machine virtuelle brute qui jouera le rôle de la VM fournie plus tard par l'équipe Système.

\subsection{Mise en place}
\begin{enumerate}
    \item \textbf{VPC} : créer un VPC (par exemple \nolinkurl{10.0.0.0/16}) avec un sous-réseau public.
    \item \textbf{Security Group} : autoriser le port 22 (SSH, limité à votre IP), les ports 80 et 443 (HTTP/HTTPS, ouverts à tous) et le port 6443 (API Kubernetes, limité à votre IP, seulement si vous pilotez le cluster depuis votre PC).
    \item \textbf{Port 5000 volontairement fermé} : le registry n'est pas exposé à Internet. Les images sont donc construites, étiquetées et poussées \textbf{depuis la VM}, ou copiées sur la VM (par exemple avec \texttt{docker save} puis \texttt{scp}).
    \item \textbf{Instance EC2} :
    \begin{itemize}
        \item AMI : Ubuntu 24.04 LTS ;
        \item type : \texttt{t3.large} (2 vCPU, 8 Go de RAM). Les instances t3 sont « burstables » : lors de tests de charge, les crédits CPU peuvent limiter les performances ;
        \item stockage : volume racine de 60 Go (gp3) ;
        \item IP : associer une \textbf{Elastic IP} (adresse fixe indispensable pour le DNS).
    \end{itemize}
\end{enumerate}

\begin{tcolorbox}[colback=yellow!10!white,colframe=orange!75!black,title=Coût]
L'Elastic IP et le volume gp3 sont facturés même instance arrêtée. Pensez à tout supprimer après la démonstration.
\end{tcolorbox}

% ==========================================
\section{Phase 2 : Installation de Docker, K3s et du registry}

\subsection{But}
Installer Docker (pour construire les images), un cluster Kubernetes léger identique dans son fonctionnement à celui du futur datacenter, et le registry privé, le tout en une seule commande.

\subsection{Utilisation}
Copier \texttt{install.sh} sur la VM, puis :
\begin{lstlisting}[language=bash]
chmod +x install.sh
./install.sh <ELASTIC_IP>
\end{lstlisting}

\textbf{Docker sans sudo.} Le script ajoute l'utilisateur au groupe \texttt{docker}, mais ce groupe ne s'applique qu'après reconnexion. Dans la session courante, utiliser \texttt{sudo docker} ou ouvrir un nouveau shell avec \texttt{newgrp docker}. Sans cela, les commandes \texttt{docker push} de la Phase~3 échouent.

\subsection{Contenu de \texttt{install.sh}}
\begin{lstlisting}[language=bash]
#!/usr/bin/env bash
# install.sh <IP_PUBLIQUE>
# Installation rejouable de la base GANDAL PaaS : Docker, K3s, registry prive.
# Teste pour Ubuntu 24.04. Autonome : ne depend d'aucun autre fichier.
set -euo pipefail

PUBLIC_IP="${1:?usage: ./install.sh <IP_PUBLIQUE>}"

echo "==> Paquets de base"
sudo apt-get update
sudo apt-get install -y curl git docker.io python3-venv
sudo systemctl enable --now docker
# Le groupe docker ne s'applique qu'apres reconnexion : utiliser 'sudo docker'
# ou 'newgrp docker' dans la session courante.
sudo usermod -aG docker "$USER"

echo "==> Configuration du registry pour K3s (HTTP sur localhost:5000)"
sudo mkdir -p /etc/rancher/k3s
sudo tee /etc/rancher/k3s/registries.yaml >/dev/null <<'EOF'
mirrors:
  "localhost:5000":
    endpoint:
      - "http://localhost:5000"
EOF

echo "==> Installation de K3s"
curl -sfL https://get.k3s.io | INSTALL_K3S_EXEC="--tls-san ${PUBLIC_IP}" sh -

mkdir -p "$HOME/.kube"
sudo cp /etc/rancher/k3s/k3s.yaml "$HOME/.kube/config"
sudo chown "$USER":"$USER" "$HOME/.kube/config"
export KUBECONFIG="$HOME/.kube/config"
grep -q KUBECONFIG "$HOME/.bashrc" || \
  echo 'export KUBECONFIG=$HOME/.kube/config' >> "$HOME/.bashrc"

# K3s peut mettre quelques secondes avant de declarer le noeud
until kubectl get nodes >/dev/null 2>&1; do sleep 3; done
kubectl wait --for=condition=Ready nodes --all --timeout=180s

echo "==> Deploiement du registry prive"
kubectl apply -f - <<'EOF'
apiVersion: v1
kind: Namespace
metadata:
  name: gandal-system
---
apiVersion: v1
kind: PersistentVolumeClaim
metadata:
  name: registry-pvc
  namespace: gandal-system
spec:
  accessModes: ["ReadWriteOnce"]
  resources:
    requests:
      storage: 10Gi
---
apiVersion: apps/v1
kind: Deployment
metadata:
  name: private-registry
  namespace: gandal-system
spec:
  replicas: 1
  strategy:
    type: Recreate
  selector:
    matchLabels:
      app: registry
  template:
    metadata:
      labels:
        app: registry
    spec:
      containers:
      - name: registry
        image: registry:2
        ports:
        - containerPort: 5000
          hostPort: 5000
        volumeMounts:
        - name: registry-storage
          mountPath: /var/lib/registry
      volumes:
      - name: registry-storage
        persistentVolumeClaim:
          claimName: registry-pvc
EOF
kubectl -n gandal-system rollout status deploy/private-registry --timeout=180s

echo "==> Environnement Python pour le script PaaS"
python3 -m venv "$HOME/gandal-venv"
"$HOME/gandal-venv/bin/pip" install --quiet kubernetes

echo
echo "Installation terminee."
echo "Reconnectez-vous (ou 'newgrp docker') pour utiliser docker sans sudo."
echo "Test registry : curl http://localhost:5000/v2/_catalog"
echo "Script PaaS   : source ~/gandal-venv/bin/activate && python gandal_deploy.py --help"
\end{lstlisting}

\subsection{Vérification}
\begin{lstlisting}[language=bash]
kubectl get nodes          # le noeud doit etre Ready
kubectl get pods -A        # traefik, coredns, metrics-server en Running
\end{lstlisting}

\subsection{Remarque sur l'adresse IP du certificat}
L'option \texttt{--tls-san} enregistre l'adresse IP dans le certificat de l'API Kubernetes. Elle ne concerne que l'accès à distance avec \texttt{kubectl} depuis un autre poste. Si l'adresse change, il suffit d'ajouter la nouvelle valeur à \texttt{tls-san} (dans \nolinkurl{/etc/rancher/k3s/config.yaml}) et de redémarrer K3s ; une réinstallation n'est pas nécessaire.

% ==========================================
\section{Phase 3 : Registry d'images privé}

\subsection{But}
Héberger les images Docker des utilisateurs dans le cluster, sans GitLab pour l'instant.

\subsection{Principe d'accès}
Le registry est déployé par \texttt{install.sh} (manifeste identique à \texttt{registry.yaml}). Il écoute sur le port 5000 \textbf{de la VM} grâce à \texttt{hostPort}. Ainsi, \nolinkurl{localhost:5000} fonctionne à la fois pour \texttt{docker push} (depuis la VM) et pour les téléchargements d'images effectués par K3s.

\subsection{Test}
\begin{lstlisting}[language=bash]
sudo docker pull nginx:alpine
sudo docker tag nginx:alpine localhost:5000/my-api:v1
sudo docker push localhost:5000/my-api:v1
curl http://localhost:5000/v2/_catalog     # doit lister my-api
\end{lstlisting}

\textbf{Remarques.}
\begin{itemize}
    \item L'image \texttt{nginx:alpine} écoute sur le port \textbf{80}. Les tests de ce document utilisent donc le port 80.
    \item Les images de test (\texttt{nginx:alpine}, \texttt{curlimages/curl}) sont téléchargées depuis Docker Hub. Sur un réseau isolé, elles doivent être fournies autrement (voir Phase~9).
    \item Le registry n'a ni authentification ni TLS : tout pod du cluster peut y lire ou y écrire. C'est acceptable pour la simulation, pas pour la production.
\end{itemize}

% ==========================================
\section{Phase 4 : Isolation, quotas et sécurité réseau}

\subsection{But}
Chaque projet utilisateur obtient un espace isolé et plafonné. Ces objets sont créés \textbf{automatiquement par \texttt{gandal\_deploy.py}} à la première création d'application d'un projet ; ils ne sont pas appliqués à la main.

\subsection{Objets créés pour chaque projet}
\begin{center}
\begin{tabular}{lll}
\toprule
\textbf{Objet} & \textbf{Nom} & \textbf{Rôle} \\
\midrule
Namespace & \nolinkurl{project-<user>-<projet>} & Unité d'isolation et de propriété \\
LimitRange & \texttt{default-limits} & Valeurs CPU/RAM par défaut des conteneurs \\
ResourceQuota & \texttt{project-quota} & Plafond global du projet (4 pods par défaut) \\
NetworkPolicy & \texttt{isolate-ingress} & Bloque le trafic entrant venant d'autres projets \\
\bottomrule
\end{tabular}
\end{center}

\subsection{Règle de la NetworkPolicy}
Sont autorisés en entrée : (1) les pods du \textbf{même namespace}, afin qu'une application puisse parler à sa base de données ; (2) le namespace \texttt{kube-system}, où tourne Traefik, afin que le trafic web atteigne l'application. Tout le reste est refusé.

\begin{lstlisting}[language=yaml]
# Resultat applique par le script (reference)
apiVersion: networking.k8s.io/v1
kind: NetworkPolicy
metadata:
  name: isolate-ingress
spec:
  podSelector: {}
  policyTypes: ["Ingress"]
  ingress:
  - from:
    - podSelector: {}
  - from:
    - namespaceSelector:
        matchLabels:
          kubernetes.io/metadata.name: kube-system
\end{lstlisting}

\textbf{Application des règles.} K3s embarque un contrôleur de NetworkPolicy (kube-router) : l'isolation fonctionne par défaut. Cela cesse d'être vrai si un autre CNI est installé ou si ce contrôleur est désactivé ; le test d'isolation du document de cadrage le vérifie en quelques secondes.

\subsection{Mapping utilisateur $\leftrightarrow$ conteneur}
Toute ressource créée porte les labels suivants :
\begin{center}
\begin{tabular}{ll}
\toprule
\textbf{Label} & \textbf{Signification} \\
\midrule
\texttt{gandal.io/owner} & Utilisateur propriétaire \\
\texttt{gandal.io/project} & Projet \\
\texttt{gandal.io/app} & Application \\
\bottomrule
\end{tabular}
\end{center}
Ces labels sont présents sur le namespace, le Deployment, les \textbf{pods}, le Service et l'Ingress. Ils permettent de retrouver tous les pods (et leurs adresses IP) d'un utilisateur par un simple sélecteur.

% ==========================================
\section{Phase 5 : Convention de nommage}

Une convention unique est utilisée partout (script, tests, documentation). Les valeurs \nolinkurl{<user>}, \nolinkurl{<projet>} et \nolinkurl{<app>} doivent respecter le format DNS-1123 (minuscules, chiffres, tirets) et faire 20 caractères maximum.

\begin{center}
\begin{tabular}{ll}
\toprule
\textbf{Ressource} & \textbf{Nom} \\
\midrule
Namespace & \nolinkurl{project-<user>-<projet>} \\
Deployment, Service, Ingress & \nolinkurl{<app>} \\
Nom DNS externe & \nolinkurl{<app>-<projet>-<user>.<domaine>} \\
Nom DNS interne au cluster & \nolinkurl{<app>.project-<user>-<projet>.svc.cluster.local} \\
\bottomrule
\end{tabular}
\end{center}

Exemple : l'application \texttt{my-api} du projet \texttt{memoire} de l'utilisateur \texttt{alice} est servie sur \nolinkurl{my-api-memoire-alice.gandal.test}.

% ==========================================
\section{Phase 6 : DNS et exposition}

\subsection{But}
Rendre chaque application accessible par un nom de domaine, sans Route 53.

\subsection{Principe}
Un seul enregistrement \textbf{wildcard} pointe vers l'IP de la VM : \nolinkurl{*.gandal.test}. Le script crée ensuite un \textbf{Ingress} par application ; Traefik route selon le nom d'hôte. Le DNS n'est donc configuré qu'une seule fois, quel que soit le nombre d'applications.

Le fichier \texttt{/etc/hosts} ne supporte pas les wildcards : il n'est pas utilisé ici.

\subsection{Option A : sslip.io (zéro configuration)}
Définir le domaine de base sur \nolinkurl{<ELASTIC_IP>.sslip.io}. Tout nom de la forme \nolinkurl{x.<ELASTIC_IP>.sslip.io} se résout automatiquement vers l'IP indiquée.
\begin{lstlisting}[language=bash]
export GANDAL_BASE_DOMAIN="<ELASTIC_IP>.sslip.io"
\end{lstlisting}

\subsection{Option B : dnsmasq sur votre PC (domaine \texttt{gandal.test})}
Sur un PC Ubuntu utilisant NetworkManager :
\begin{lstlisting}[language=bash]
sudo apt install -y dnsmasq-base
sudo sed -i 's/^\[main\]/[main]\ndns=dnsmasq/' /etc/NetworkManager/NetworkManager.conf
echo "address=/gandal.test/<ELASTIC_IP>" | \
  sudo tee /etc/NetworkManager/dnsmasq.d/gandal.conf
sudo systemctl restart NetworkManager
getent hosts test.gandal.test       # doit renvoyer l'Elastic IP
\end{lstlisting}

\subsection{Remarque}
CoreDNS est le DNS \textbf{interne} du cluster (résolution des Services entre pods). Il n'intervient pas dans l'accès externe des clients, qui repose sur le wildcard ci-dessus.

% ==========================================
\section{Phase 7 : Automatisation, le cœur du PaaS}

\subsection{But}
Le script \texttt{gandal\_deploy.py} fournit quatre commandes :
\begin{itemize}
    \item \texttt{deploy} : crée ou met à jour le namespace, les quotas, l'isolation, le Deployment, le Service et l'Ingress ;
    \item \texttt{list} : liste les pods avec leur propriétaire, leur statut et leur \textbf{adresse IP} ;
    \item \texttt{status} : état d'une application (réplicas prêts, URL, IP des pods) ;
    \item \texttt{delete} : supprime une application, ou tout le projet avec \texttt{--purge}.
\end{itemize}

Le script est \textbf{idempotent} : relancer \texttt{deploy} met à jour la ressource existante (nouvelle image, nouveau nombre de réplicas) au lieu de l'ignorer. Les valeurs dépendantes de l'environnement sont lues dans les variables \nolinkurl{GANDAL_BASE_DOMAIN}, \nolinkurl{GANDAL_INGRESS_CLASS}, \nolinkurl{GANDAL_REPLICAS} et \nolinkurl{GANDAL_QUOTA_*}.

\subsection{Utilisation}
\begin{lstlisting}[language=bash]
source ~/gandal-venv/bin/activate
python gandal_deploy.py deploy alice memoire my-api localhost:5000/my-api:v1 80
python gandal_deploy.py list --user alice
python gandal_deploy.py status alice memoire my-api
python gandal_deploy.py delete alice memoire my-api
python gandal_deploy.py delete alice memoire --purge
\end{lstlisting}

Le script s'exécute sur la VM. Pour le lancer depuis votre PC, copier \texttt{\textasciitilde/.kube/config}, remplacer \nolinkurl{127.0.0.1} par l'Elastic IP dans le champ \texttt{server}, et ouvrir le port 6443 à votre IP.

\subsection{Code}
\begin{lstlisting}[language=Python]
#!/usr/bin/env python3
"""GANDAL PaaS (simulation) : deploy / list / status / delete."""
import argparse
import os
import re
import sys

from kubernetes import client, config
from kubernetes.client.rest import ApiException

BASE_DOMAIN = os.getenv("GANDAL_BASE_DOMAIN", "gandal.test")
INGRESS_CLASS = os.getenv("GANDAL_INGRESS_CLASS", "traefik")
REPLICAS = int(os.getenv("GANDAL_REPLICAS", "2"))
QUOTA = {
    "requests.cpu": os.getenv("GANDAL_QUOTA_REQ_CPU", "1"),
    "requests.memory": os.getenv("GANDAL_QUOTA_REQ_MEM", "1Gi"),
    "limits.cpu": os.getenv("GANDAL_QUOTA_LIM_CPU", "2"),
    "limits.memory": os.getenv("GANDAL_QUOTA_LIM_MEM", "2Gi"),
    "pods": os.getenv("GANDAL_QUOTA_PODS", "4"),
}
NAME_RE = re.compile(r"^[a-z0-9]([-a-z0-9]*[a-z0-9])?$")


def check(label, value):
    if not NAME_RE.match(value) or len(value) > 20:
        sys.exit(f"{label} invalide '{value}' : minuscules, chiffres, "
                 "tirets, 20 caracteres max")


def ns_of(user, project):
    return f"project-{user}-{project}"


def host_of(user, project, app):
    return f"{app}-{project}-{user}.{BASE_DOMAIN}"


def labels(user, project, app=None):
    d = {"gandal.io/owner": user, "gandal.io/project": project}
    if app:
        d["gandal.io/app"] = app
    return d


def upsert(create, patch, name, body):
    """Cree la ressource ; si elle existe deja (409), la met a jour."""
    try:
        create(body)
        print(f"  cree    : {name}")
    except ApiException as e:
        if e.status != 409:
            raise
        patch(name, body)
        print(f"  a jour  : {name}")


def deploy(a):
    for label, value in (("user", a.user), ("projet", a.project),
                         ("app", a.app)):
        check(label, value)
    config.load_kube_config()
    v1, apps, net = (client.CoreV1Api(), client.AppsV1Api(),
                     client.NetworkingV1Api())
    ns = ns_of(a.user, a.project)
    lab = labels(a.user, a.project, a.app)
    print(f"Deploiement de {a.app} dans {ns}")

    # 1. Namespace (porte owner et project)
    upsert(v1.create_namespace, lambda n, b: v1.patch_namespace(n, b), ns,
           {"metadata": {"name": ns, "labels": labels(a.user, a.project)}})

    # 2. LimitRange : valeurs par defaut des conteneurs
    upsert(lambda b: v1.create_namespaced_limit_range(ns, b),
           lambda n, b: v1.patch_namespaced_limit_range(n, ns, b),
           "default-limits",
           {"metadata": {"name": "default-limits"},
            "spec": {"limits": [{
                "type": "Container",
                "default": {"cpu": "500m", "memory": "512Mi"},
                "defaultRequest": {"cpu": "100m", "memory": "128Mi"}}]}})

    # 3. ResourceQuota : plafond du projet
    upsert(lambda b: v1.create_namespaced_resource_quota(ns, b),
           lambda n, b: v1.patch_namespaced_resource_quota(n, ns, b),
           "project-quota",
           {"metadata": {"name": "project-quota"}, "spec": {"hard": QUOTA}})

    # 4. NetworkPolicy : meme namespace + kube-system (Traefik)
    upsert(lambda b: net.create_namespaced_network_policy(ns, b),
           lambda n, b: net.patch_namespaced_network_policy(n, ns, b),
           "isolate-ingress",
           {"metadata": {"name": "isolate-ingress"},
            "spec": {"podSelector": {}, "policyTypes": ["Ingress"],
                     "ingress": [
                         {"from": [{"podSelector": {}}]},
                         {"from": [{"namespaceSelector": {"matchLabels": {
                             "kubernetes.io/metadata.name":
                                 "kube-system"}}}]}]}})

    # 5. Deployment (labels de propriete sur les pods)
    probe = {"tcpSocket": {"port": a.port}, "periodSeconds": 5}
    upsert(lambda b: apps.create_namespaced_deployment(ns, b),
           lambda n, b: apps.patch_namespaced_deployment(n, ns, b),
           a.app,
           {"metadata": {"name": a.app, "labels": lab},
            "spec": {
                "replicas": a.replicas,
                "selector": {"matchLabels": {"gandal.io/app": a.app}},
                "template": {
                    "metadata": {"labels": lab},
                    "spec": {"containers": [{
                        "name": a.app, "image": a.image,
                        "ports": [{"containerPort": a.port}],
                        "readinessProbe": dict(probe,
                                               initialDelaySeconds=3),
                        "livenessProbe": dict(probe,
                                              initialDelaySeconds=10),
                        "resources": {
                            "requests": {"cpu": "100m",
                                         "memory": "128Mi"},
                            "limits": {"cpu": "500m",
                                       "memory": "512Mi"}}}]}}}})

    # 6. Service (nom stable a l'interieur du cluster)
    upsert(lambda b: v1.create_namespaced_service(ns, b),
           lambda n, b: v1.patch_namespaced_service(n, ns, b),
           a.app,
           {"metadata": {"name": a.app, "labels": lab},
            "spec": {"selector": {"gandal.io/app": a.app},
                     "ports": [{"port": 80, "targetPort": a.port}]}})

    # 7. Ingress (nom DNS externe)
    host = host_of(a.user, a.project, a.app)
    upsert(lambda b: net.create_namespaced_ingress(ns, b),
           lambda n, b: net.patch_namespaced_ingress(n, ns, b),
           a.app,
           {"metadata": {"name": a.app, "labels": lab},
            "spec": {"ingressClassName": INGRESS_CLASS,
                     "rules": [{"host": host, "http": {"paths": [{
                         "path": "/", "pathType": "Prefix",
                         "backend": {"service": {
                             "name": a.app,
                             "port": {"number": 80}}}}]}}]}})

    print(f"Succes. Acces : http://{host}")


def list_pods(a):
    config.load_kube_config()
    v1 = client.CoreV1Api()
    sel = f"gandal.io/owner={a.user}" if a.user else "gandal.io/owner"
    pods = v1.list_pod_for_all_namespaces(label_selector=sel).items
    print(f"{'OWNER':12}{'PROJECT':14}{'APP':16}{'POD':40}"
          f"{'STATUS':10}IP")
    for p in pods:
        lb = p.metadata.labels or {}
        print(f"{lb.get('gandal.io/owner', ''):12}"
              f"{lb.get('gandal.io/project', ''):14}"
              f"{lb.get('gandal.io/app', ''):16}"
              f"{p.metadata.name:40}{p.status.phase:10}"
              f"{p.status.pod_ip or '-'}")


def status(a):
    config.load_kube_config()
    v1, apps = client.CoreV1Api(), client.AppsV1Api()
    ns = ns_of(a.user, a.project)
    d = apps.read_namespaced_deployment(a.app, ns)
    print(f"Replicas prets : {d.status.ready_replicas or 0}/{d.spec.replicas}")
    print(f"URL            : http://{host_of(a.user, a.project, a.app)}")
    pods = v1.list_namespaced_pod(
        ns, label_selector=f"gandal.io/app={a.app}").items
    for p in pods:
        print(f"Pod {p.metadata.name} : {p.status.phase}, "
              f"IP {p.status.pod_ip}")


def delete(a):
    config.load_kube_config()
    v1, apps, net = (client.CoreV1Api(), client.AppsV1Api(),
                     client.NetworkingV1Api())
    ns = ns_of(a.user, a.project)
    if a.purge:
        v1.delete_namespace(ns)
        print(f"Projet {ns} supprime")
        return
    if not a.app:
        sys.exit("indiquer l'application, ou utiliser --purge")
    for fn in (lambda: net.delete_namespaced_ingress(a.app, ns),
               lambda: v1.delete_namespaced_service(a.app, ns),
               lambda: apps.delete_namespaced_deployment(a.app, ns)):
        try:
            fn()
        except ApiException as e:
            if e.status != 404:
                raise
    print(f"Application {a.app} supprimee")


def main():
    p = argparse.ArgumentParser(description="GANDAL PaaS (simulation)")
    sub = p.add_subparsers(dest="cmd", required=True)

    d = sub.add_parser("deploy", help="creer ou mettre a jour une application")
    for n in ("user", "project", "app", "image"):
        d.add_argument(n)
    d.add_argument("port", type=int)
    d.add_argument("--replicas", type=int, default=REPLICAS)
    d.set_defaults(fn=deploy)

    ls = sub.add_parser("list", help="lister les pods avec proprietaire et IP")
    ls.add_argument("--user")
    ls.set_defaults(fn=list_pods)

    s = sub.add_parser("status", help="etat d'une application")
    for n in ("user", "project", "app"):
        s.add_argument(n)
    s.set_defaults(fn=status)

    x = sub.add_parser("delete", help="supprimer une application ou un projet")
    x.add_argument("user")
    x.add_argument("project")
    x.add_argument("app", nargs="?")
    x.add_argument("--purge", action="store_true")
    x.set_defaults(fn=delete)

    args = p.parse_args()
    args.fn(args)


if __name__ == "__main__":
    main()
\end{lstlisting}

% ==========================================
\section{Phase 8 : Validation et résilience}

Les tests détaillés figurent dans le document de cadrage. Les commandes de base :

\begin{lstlisting}[language=bash]
# Self-healing : supprimer un pod, Kubernetes le recree
kubectl delete pod -n project-alice-memoire -l gandal.io/app=my-api
kubectl get pods -n project-alice-memoire -w

# Metriques (metrics-server est inclus dans K3s)
kubectl top pods -n project-alice-memoire

# Logs applicatifs
kubectl logs -n project-alice-memoire -l gandal.io/app=my-api

# IP des pods d'un utilisateur
python gandal_deploy.py list --user alice
\end{lstlisting}

% ==========================================
\section{Phase 9 : Passage à GANDAL (XCP-ng)}

\subsection{Ce qui ne change pas}
Le manifeste du registry, la logique du script Python, les Namespaces, quotas, NetworkPolicies, labels et Ingress, ainsi que le script \texttt{install.sh}.

\subsection{Ce qui change}
\begin{center}
\begin{tabular}{p{4cm}p{4.4cm}p{5.2cm}}
\toprule
\textbf{Élément} & \textbf{AWS (simulation)} & \textbf{GANDAL (XCP-ng)} \\
\midrule
Adresse d'entrée & Elastic IP (NAT AWS) & IP de la VM fournie par Système/Réseau \\
Domaine & \nolinkurl{gandal.test} ou sslip.io & Domaine GANDAL, wildcard géré par l'équipe Réseau \\
Variable à modifier & \nolinkurl{GANDAL_BASE_DOMAIN} & idem \\
TLS (HTTPS) & Absent & cert-manager ou certificats fournis \\
Registry & HTTP sans authentification & TLS et authentification (ou registry GitLab plus tard) \\
Pare-feu & Security Group & Règles Réseau/Sécurité \\
Accès Internet sortant & Libre & À vérifier : \texttt{install.sh} télécharge K3s, des paquets et des images \\
CNI et NetworkPolicies & kube-router intégré à K3s & Valable tant qu'aucun autre CNI n'est installé ; sinon vérifier \\
Adresse IP du certificat & \texttt{--tls-san} avec l'Elastic IP & Idem avec l'IP de la VM ; ajouter toute nouvelle IP puis redémarrer K3s \\
\bottomrule
\end{tabular}
\end{center}

\subsection{Si le réseau GANDAL est isolé d'Internet}
Les téléchargements de \texttt{install.sh} (paquets apt, script et binaire K3s, image \texttt{registry:2}) et des tests (\texttt{nginx:alpine}, \texttt{curlimages/curl}) échoueront. Dans ce cas :
\begin{itemize}
    \item demander à Système/Réseau quelles sorties sont autorisées (ou un miroir interne) ;
    \item préparer les images sur une machine connectée avec \texttt{docker save}, les copier sur la VM, les charger avec \texttt{docker load}, puis les pousser dans le registry local ;
    \item utiliser l'installation hors ligne de K3s (binaire et images fournis à l'avance).
\end{itemize}

\subsection{Procédure}
\begin{enumerate}
    \item Recevoir la VM et son IP.
    \item Vérifier l'accès Internet sortant (ou préparer l'installation hors ligne).
    \item Exécuter \texttt{install.sh} avec cette IP.
    \item Définir \nolinkurl{GANDAL_BASE_DOMAIN} et faire créer le wildcard DNS par l'équipe Réseau.
    \item Relancer les tests de validation du document de cadrage.
\end{enumerate}

% ==========================================
\section{Conclusion}

Cette procédure donne à la cellule PaaS un prototype conforme aux exigences : déploiements par conteneurs, orchestration par Kubernetes, association utilisateur--pods par labels, récupération des adresses IP des pods, et nom DNS d'accès par application. La migration vers XCP-ng se limite à relancer \texttt{install.sh} et à ajuster les variables et points listés à la Phase~9.

\end{document}
